\section{What this framework changes}

The repository's many proof packets already exhibit substantial common
structure: rational approximations, finite algebra, explicit inequalities,
primitive recursive algorithms, and an audit of the allowed induction.
They also repeatedly encounter the same missing interfaces. A geometric
argument may produce a finiteness assertion without a coordinate bound.
An analytic argument may require an exact zero test on a Cauchy name.
A Lean argument may prove a statement about natural numbers while leaving
the object-language arithmetic derivation unconstructed. The useful next
step is to make these interfaces reusable and then prove new instances.

\subsection{The proposed unit of mathematical development}

An observation request is a finite datum: a precision, a polynomial degree,
a bounded list of points, a finite matrix, a requested range of indices,
or a finite assertion about one of these. A construction should specify:
\begin{enumerate}
\item a fixed primitive recursive representation of its inputs and outputs;
\item the finite observations of the input used for this request;
\item an explicit primitive recursive output procedure;
\item a finite correctness predicate, including all necessary promises;
\item a proof whose loop invariants become quantifier-free arithmetic.
\end{enumerate}
Theorem~\ref{thm:finite-demand-transfer} supplies the general transfer step.
It allows adaptive input demands and global hypotheses, provided the
induction itself is performed at a fixed finite request. The global
interpretation is recovered afterwards by universal introduction.

This permits quite different objects to share one proof discipline. An
exact rational algebraic object uses finite equality tests. A real or
complex name uses finite coherence and error tests. A function uses an
evaluator together with a quantitative continuity contract. A finite
geometric configuration uses incidence and nondegeneracy certificates.
A count bound uses a contradiction for a supplied list of too many
distinct points. The framework does not collapse these representations
into one stronger existence claim.

\subsection{The main mathematical outcomes}

The principal positive results are the following.
\begin{enumerate}
\item \textbf{A direct arithmetic transfer theorem.}
Bounded truth, bounded witness selection, adaptive finite demand, and
composition are proved explicitly. A restricted proof discipline has an
effective syntactic compilation into iPRA. Its primitive contracts must
already be proved; the theorem does not supply missing mathematics.

\item \textbf{A separated simple branch for Cauchy coefficients.}
Theorem~\ref{thm:simple-branch} replaces exact algebraic coefficient
codes by a different sufficient interface: arbitrary fixed PR complex
coefficient names, a supplied root at $X=0$, and an explicit separation
of the derivative from zero. Finite polynomial iteration produces
individually accessible coefficient names and quantitative convergence.
The derivative witness is the exact point where the singular Puiseux
obstruction is excluded.

\item \textbf{Common finite windows without a global modulus.}
Theorem~\ref{thm:common-window} combines finite variation bounds and
finite exceptional-set bounds by addition of budgets. It gives one
bounded search that satisfies all observations on the same window.
It includes nonmonotone sequences. A halting example proves why this
output cannot generally be promoted to a Cauchy modulus.

\item \textbf{A concrete repair of a Diophantine geometric dependency.}
Theorem~\ref{thm:hyperelliptic-fiber} supplies a sharp finite zero-count
lemma on odd-degree hyperelliptic equations. It replaces a divisor-based
argument in a specified part of the Siegel program. Conditional on the
still-needed finite covering theorem, it gives the point-count bound
$ML$ for the explicit evaluation vectors used in
Corollary~\ref{cor:hyperelliptic-cardinality}.

\item \textbf{A general finite algebra interface.}
Theorem~\ref{thm:finite-algebra-zeros} converts explicit polynomial
spanning identities into a bound on the number of rational solutions.
The certificate can be checked by coefficient arithmetic. Its proof
uses separator polynomials and finite row reduction, and does not
assume a generic algebraic-geometry decision procedure.
\end{enumerate}

The elementary ingredients of these results are established mathematics.
The contribution of this report is their integration into a common iPRA
discipline, the explicit arithmetic and precision details, and their
application to the specific unresolved contracts in the repository.
No claim of priority for a new classical analytic or algebraic theorem
is intended.

\subsection{Three distinctions that preserve the intended theorem}

\paragraph{Counting points versus finding all points.}
A proof that every list of $C+1$ distinct solutions is impossible gives
a cardinality bound. It does not tell a bounded search where the last
solution lies. For an arbitrary decidable predicate, even a bound of
one permits a sole solution at an arbitrarily late index. A global
height bound is additional information. The finite algebra and
hyperelliptic results in this report are count results.

\paragraph{Proving a finite certificate sound versus producing one.}
A coefficient identity can be verified and its consequence can be
proved uniformly in a supplied certificate. This does not show that
every input in a proposed theorem admits such a certificate. For a
uniform positive result, either supply the certificate as input or
prove a primitive recursive constructor and its correctness. This
distinction applies to algebraic splitting, exceptional subspace
covers, root isolation, and proof traces alike.

\paragraph{Native Lean truth versus an embedded arithmetic proof.}
Lean is the metatheory in the intended formalization. The target theory
is represented by its syntax and derivation type. A theorem about
\texttt{Nat}, \texttt{Rat}, or \texttt{Complex} in Lean is useful
mathematics, but the required final endpoint is a value of the repository
type \texttt{Prf PRAAx [] formula}. The archive's formalization plan
keeps these obligations separate, including which files were actually
compiled in this work.

\subsection{How to read and use the report}

Section~\ref{sec:finite-demand-framework} gives the common foundations.
Section~\ref{sec:simplebranch} develops the analytic repair and a finite
mesh method. Section~\ref{sec:stability} gives the stability selector.
Section~\ref{sec:repository-repair} audits the hard packets and develops
the geometric and finite algebra repairs. The final sections specify
the exact Lean interfaces, remaining obligations, reproducible checks,
and the boundaries of the claims.

The manuscript is intended to stand on its own. The repository is used
to identify the target foundations and the particular gaps; no unverified
analytic theorem from another packet is a hidden premise of the new
finite polynomial, stability, or zero-count proofs. The conditional
covering corollary deliberately leaves its covering theorem as an
explicit hypothesis.
